\documentclass[a4paper,11pt]{book}
%\documentclass[a4paper,twoside,11pt,titlepage]{book}
\usepackage{listings}
\usepackage[utf8]{inputenc}
\usepackage[spanish]{babel}

\usepackage{glossaries}
%\usepackage{titlesec}
%\usepackage{pailatino}
% Para el uso de FloatBarrier
\usepackage{placeins}

% Para ajustar 
\usepackage{adjustbox}



% Para la representación de tablas
\usepackage{tabularx}
\usepackage{multirow}
\usepackage{booktabs}
\usepackage[table,xcdraw]{xcolor}

\newcommand{\mc}[3]{\multicolumn{#1}{#2}{#3}}
\newcommand{\mr}[3]{\multirow{#1}{#2}{#3}}
\newcommand{\nl}{\tabularnewline}

% Para la realización de digramas de gantt
\usepackage{pgfgantt}
\usepackage[spanish]{datetime2}
\usepackage{pgfcalendar}

\renewcommand*{\pgfcalendarmonthname}[1]{%
  \ifcase#1\or
  enero\or febrero\or marzo\or abril\or mayo\or junio\or
  julio\or agosto\or septiembre\or octubre\or noviembre\or diciembre\fi
}

% Para la representación de esquemas
\usepackage{tikz}
\usetikzlibrary{positioning, arrows}

 % PARA CODIGO PYTHON
\usepackage{listings}

\lstdefinestyle{mybash}{
    language=bash,
    breaklines=true,
    basicstyle=\ttfamily\small,
    keywordstyle=\color{blue},
    stringstyle=\color{red},
    commentstyle=\color{olive},
    morecomment=[l]{\#},
    morekeywords={screen, ifconfig, ovs-vsctl, sshpass, ssh, ovsdb-server, ovs-vswitchd},
    emph={then, do, done, fi, for, if, in},
    emphstyle={\color{blue}},
    identifierstyle=\color{black},
    numberstyle=\tiny\color{gray},
    frame=single,
    backgroundcolor=\color{white}
}

\lstdefinestyle{mypy}{
    language=Python,
    breaklines=true,
    basicstyle=\ttfamily,
    keywordstyle=\color{blue},
    stringstyle=\color{red},
    commentstyle=\color{green},
    morecomment=[l]{\#},
}


\decimalpoint
\usepackage{dcolumn}
\newcolumntype{.}{D{.}{\esperiod}{-1}}
\makeatletter
\addto\shorthandsspanish{\let\esperiod\es@period@code}
\makeatother


%\usepackage[chapter]{algorithm}
\RequirePackage{verbatim}
%\RequirePackage[Glenn]{fncychap}
\usepackage{fancyhdr}
\usepackage{graphicx}
\usepackage{afterpage}

\usepackage{longtable}

\usepackage[pdfborder={000}]{hyperref} %referencia

% ********************************************************************
% Re-usable information
% ********************************************************************
\newcommand{\myTitle}{Título del proyecto\xspace}
\newcommand{\myDegree}{Grado en ...\xspace}
\newcommand{\myName}{Nombre Apllido1 Apellido2 (alumno)\xspace}
\newcommand{\myProf}{Nombre Apllido1 Apellido2 (tutor1)\xspace}
\newcommand{\myOtherProf}{Nombre Apllido1 Apellido2 (tutor2)\xspace}
%\newcommand{\mySupervisor}{Put name here\xspace}
\newcommand{\myFaculty}{Escuela Técnica Superior de Ingenierías Informática y de
Telecomunicación\xspace}
\newcommand{\myFacultyShort}{E.T.S. de Ingenierías Informática y de
Telecomunicación\xspace}
\newcommand{\myDepartment}{Departamento de ...\xspace}
\newcommand{\myUni}{\protect{Universidad de Granada}\xspace}
\newcommand{\myLocation}{Granada\xspace}
\newcommand{\myTime}{\today\xspace}
\newcommand{\myVersion}{Version 0.1\xspace}


\addto\captionsspanish{
  \renewcommand{\listtablename}{Índice de tablas}
  \renewcommand{\tablename}{Tabla}
}


\hypersetup{
pdfauthor = {\myName (email (en) ugr (punto) es)},
pdftitle = {\myTitle},
pdfsubject = {},
pdfkeywords = {palabra_clave1, palabra_clave2, palabra_clave3, ...},
pdfcreator = {LaTeX con el paquete ....},
pdfproducer = {pdflatex}
}

%\hyphenation{}


%\usepackage{doxygen/doxygen}
%\usepackage{pdfpages}
\usepackage{url}
\usepackage{colortbl,longtable}
\usepackage[stable]{footmisc}
%\usepackage{index}

\renewcommand{\tablename}{Tabla}
%\usepackage[style=long, cols=2,border=plain,toc=true,number=none]{glossary}
% \makeglossary

% Definición de comandos que me son tiles:
%\renewcommand{\indexname}{Índice alfabético}
%\renewcommand{\glossaryname}{Glosario}

\pagestyle{fancy}
\fancyhf{}
\fancyhead[LO]{\leftmark}
\fancyhead[RE]{\rightmark}
\fancyhead[RO,LE]{\textbf{\thepage}}
\renewcommand{\chaptermark}[1]{\markboth{\textbf{#1}}{}}
\renewcommand{\sectionmark}[1]{\markright{\textbf{\thesection. #1}}}

\setlength{\headheight}{1.5\headheight}

\newcommand{\HRule}{\rule{\linewidth}{0.5mm}}
%Definimos los tipos teorema, ejemplo y definición podremos usar estos tipos
%simplemente poniendo \begin{teorema} \end{teorema} ...
\newtheorem{teorema}{Teorema}[chapter]
\newtheorem{ejemplo}{Ejemplo}[chapter]
\newtheorem{definicion}{Definición}[chapter]

\definecolor{gray97}{gray}{.97}
\definecolor{gray75}{gray}{.75}
\definecolor{gray45}{gray}{.45}
\definecolor{gray30}{gray}{.94}

\lstset{ frame=Ltb,
     framerule=0.5pt,
     aboveskip=0.5cm,
     framextopmargin=3pt,
     framexbottommargin=3pt,
     framexleftmargin=0.1cm,
     framesep=0pt,
     rulesep=.4pt,
     backgroundcolor=\color{gray97},
     rulesepcolor=\color{black},
     %
     stringstyle=\ttfamily,
     showstringspaces = false,
     basicstyle=\scriptsize\ttfamily,
     commentstyle=\color{gray45},
     keywordstyle=\bfseries,
     %
     numbers=left,
     numbersep=6pt,
     numberstyle=\tiny,
     numberfirstline = false,
     breaklines=true,
   }
 
% minimizar fragmentado de listados
\lstnewenvironment{listing}[1][]
   {\lstset{#1}\pagebreak[0]}{\pagebreak[0]}

\lstdefinestyle{CodigoC}
   {
	basicstyle=\scriptsize,
	frame=single,
	language=C,
	numbers=left
   }
\lstdefinestyle{CodigoC++}
   {
	basicstyle=\small,
	frame=single,
	backgroundcolor=\color{gray30},
	language=C++,
	numbers=left
   }

 
\lstdefinestyle{Consola}
   {basicstyle=\scriptsize\bf\ttfamily,
    backgroundcolor=\color{gray30},
    frame=single,
    numbers=none
   }


\newcommand{\bigrule}{\titlerule[0.5mm]}





%Para conseguir que en las páginas en blanco no ponga cabecerass
\makeatletter
\def\clearpage{%
  \ifvmode
    \ifnum \@dbltopnum =\m@ne
      \ifdim \pagetotal <\topskip
        \hbox{}
      \fi
    \fi
  \fi
  \newpage
  \thispagestyle{empty}
  \write\m@ne{}
  \vbox{}
  \penalty -\@Mi
}
\makeatother

\usepackage{pdfpages}
\usepackage{appendix}
\usepackage{float}

\begin{document}

    % RODGOM - INCLUYO ESTO PARA LA NUMERACIÓN CORRECTA
    \frontmatter
    % Plantilla portada UGR
    \begin{titlepage}
 
 
\newlength{\centeroffset}
\setlength{\centeroffset}{-0.5\oddsidemargin}
\addtolength{\centeroffset}{0.5\evensidemargin}
\thispagestyle{empty}

\noindent\hspace*{\centeroffset}\begin{minipage}{\textwidth}

\centering
\includegraphics[width=0.9\textwidth]{imagenes/logo_ugr.jpg}\\[1.4cm]

\textsc{ \Large TRABAJO FIN DE GRADO\\[0.2cm]}
\textsc{GRADO EN INGENIERÍA INFORMÁTICA}\\[1cm]
% Upper part of the page
% 
% Title
{\Huge\bfseries  Sistema de Tutorización Inteligente para la Enseñanza de Programación a Niños\\[0.2cm]}
\noindent\rule[-1ex]{\textwidth}{3pt}\\[3.5ex]
{\large\bfseries Diseño e Implementación de una Plataforma Web con un Sistema de Tutorización Inteligente para facilitar la Enseñanza}
\end{minipage}

\vspace{0.5cm}
\noindent\hspace*{\centeroffset}\begin{minipage}{\textwidth}
\centering

\textbf{Autor}\\ {Noelia Carrasco Vilar}\\[1cm]
\textbf{Director}\\ {José Manuel Benítez Sánchez}\\[1cm]
\includegraphics[width=0.3\textwidth]{imagenes/etsiit_logo.png}\\[0.5cm]
\textsc{Escuela Técnica Superior de Ingenierías Informática y de Telecomunicación}\\
\textsc{---}\\
Granada, Noviembre de 2023
\end{minipage}
\end{titlepage}

    
    

    % Plantilla prefacio UGR
    \chapter*{}
\thispagestyle{empty}
\begin{center}

\large\bfseries Sistema de Tutorización Inteligente para la Enseñanza de Programación a Niños\\[1cm]

\end{center}

\begin{center}

Noelia Carrasco Vilar \\

\end{center}

\vspace{0.5cm}
\noindent{\textbf{Palabras clave}:
\textit{Sistema de Tutorización Inteligente},
\textit{Educación Infantil},
\textit{Análisis},
\textit{Automatización},
\textit{Programación}}
\vspace{0.5cm}

\noindent{\textbf{Resumen}}\\

La enseñanza de la programación en edades tempranas se ha vuelto fundamental en la sociedad actual. Es en este contexto donde los Sistemas de Tutorización Inteligente (ITS) emergen como herramientas potenciales para apoyar un aprendizaje más personalizado y adaptado a las necesidades individuales de cada estudiante.

Este trabajo presenta el diseño, desarrollo y evaluación de un ITS enfocado en la enseñanza de programación. El sistema propuesto aplica una automatización en la corrección de los ejercicios con tal de poder ofrecer una buena retroalimentación en tiempo real. Además de adaptarse a las necesidades individuales de los estudiantes y monitorizar su progreso. 

El proceso de desarrollo del ITS se llevó a cabo en diversas fases. En primer lugar, se establecieron los fundamentos pedagógicos y se identificaron las necesidades específicas del público infantil. Posteriormente, se integraron técnicas avanzadas de corrección y elección de ejercicios con tal de garantizar la adaptabilidad y efectividad del sistema.

Para evaluar la eficacia, se realizó una prueba con un pequeño grupo de estudiantes, en colaboración con la empresa CODELEARN S.L. Los resultados iniciales sugirieron una posible mejora en la retención de conocimientos y en el rendimiento académico de los participantes, validando así la propuesta de este TFG.

En conclusión, este Trabajo de Fin de Grado aporta una solución innovadora al campo de la educación tecnológica infantil, demostrando la viabilidad y eficacia de los Sistemas de Tutorización Inteligente dentro del mundo educativo de la programación más allá de la corrección rutinaria.


% ABSTRACT

\chapter*{}
\thispagestyle{empty}

\begin{center}

{\large\bfseries Intelligent Tutoring System for Teaching Programming to Children}\\[1cm]

\end{center}

\begin{center}

Noelia Carrasco Vilar\\

\end{center}

\vspace{0.5cm}
\noindent{\textbf{Keywords}:
\textit{Intelligent Tutoring System},
\textit{Early Childhood Education},
\textit{Analysis},
\textit{Automation},
\textit{Programming}}
\vspace{0.5cm}

\noindent{\textbf{Abstract}}\\

Teaching programming at an early age has become essential in today's society. It is in this context that Intelligent Tutoring Systems (ITS) emerge as potential tools to support more personalized learning tailored to the individual needs of each student.

This work presents the design, development, and evaluation of an ITS focused on teaching programming to children. The proposed system applies automation techniques in exercise correction to provide real-time feedback on the student's understanding. It is also capable of adapting to individual student needs and monitoring their progress.

The development of the ITS was carried out in various phases. Firstly, pedagogical foundations were established, and the specific needs of the child audience were identified. Advanced correction and exercise selection techniques were then integrated to ensure the adaptability and effectiveness of the system.

To evaluate the effectiveness of the ITS, tests and experiments were conducted with a small group , in collaboration with the company CODELEARN S.L. Initial results suggest a possible improvement in knowledge retention and academic performance of the participants, thus validating the proposal of this Bachelor's Thesis.

In conclusion, this Bachelor's Thesis provides an innovative solution to the field of early childhood technological education, demonstrating the feasibility and effectiveness of Intelligent Tutoring Systems in the educational programming world.

\chapter*{}
\thispagestyle{empty}

\noindent\rule[-1ex]{\textwidth}{2pt}\\[4.5ex]

Yo, \textbf{Noelia Carrasco Vilar}, alumna de la titulación Grado en Ingeniería Informática de la \textbf{Escuela Técnica Superior
de Ingenierías Informática y de Telecomunicación de la Universidad de
Granada}, con DNI AAAAAAAAA, autorizo la
ubicación de la siguiente copia de mi Trabajo Fin de Grado en la biblioteca
del centro para que pueda ser
consultada por las personas que lo deseen.

\vspace{6cm}

\noindent Fdo: Noelia Carrasco Vilar

\vspace{2cm}

\begin{flushright}
Granada a 13 de Noviembre de 2023.
\end{flushright}


\chapter*{}
\thispagestyle{empty}

\noindent\rule[-1ex]{\textwidth}{2pt}\\[4.5ex]

D. \textbf{José Manuel Benítez Sánchez}, catedrático del departamento de Ciencias de la Computación e Inteligencia Artificial de la Universidad de Granada.


\vspace{0.5cm}

\textbf{Informa:}

\vspace{0.5cm}

Que el presente trabajo, titulado \textit{\textbf{Sistema de tutorización inteligente para la enseñanza de programación a niños}}, ha sido realizado bajo su supervisión por \textbf{Noelia Carrasco Vilar}, y autoriza la defensa de dicho trabajo ante el tribunal que corresponda.

\vspace{0.5cm}

Y para que conste, expide y firma el presente informe en Granada a 13 de
Noviembre de 2023.

\vspace{1cm}

\textbf{El director:}

\vspace{5cm}

\noindent \textbf{José Manuel Benítez Sánchez}

\chapter*{Agradecimientos}

Quiero comenzar agradeciendo a los profesores que han realizado un compromiso para la formación académica de mi promoción. Su pasión y dedicación fomentaron mi interés por la informática, lo cual me ha llevado a realizar este TFG. Un agradecimiento especial a mi tutor del TFG, quien guío y ayudó este proyecto.

Agradezco de manera especial a la empresa Codelearn S.L. por su apoyo y colaboración en este proyecto. Su experiencia y recursos han sido cruciales para el desarrollo del proyecto.

No puedo pasar por alto el papel fundamental que mi familia ha desempeñado en mi vida, tanto académica como personal. A mis padres y a mi hermano, cuya influencia ha sido determinante en mi formación. Dedico este trabajo y los proyectos futuros, como reconocimiento a todos los esfuerzos que han realizado para permitirme alcanzar mis metas académicas.  También deseo honrar la memoria de mis abuelos, quienes, presentes o ausentes, han sido una fuente de inspiración constante a lo largo de mi trayectoria.




    % Índice de contenidos
    \newpage
    \tableofcontents

    % Índice de imágenes y tablas
    \newpage
    \listoffigures

    % Si hay suficientes se incluirá dicho índice
    \listoftables
    \newpage

    % RODGOM - INCLUYO ESTO PARA LA NUMERACIÓN CORRECTA
    \mainmatter
    % Introducción y descripción del problema
    \chapter{Introducción} \label{chap:introduction}

\section{Contexto Actual y Relevancia de la Programación}
En la era actual, la digitalización está en todos los aspectos de nuestra vida. Afecta desde cómo nos comunicamos hasta como realizamos nuestras tareas cotidianas. Dentro de este contexto, la programación ha transcendido su función original como mera herramienta técnica para convertirse en una habilidad esencial del siglo XXI. No solo es importante para aquellos que buscan una carrera tecnológica, sino también para el público que desea entender y utilizar eficazmente las herramientas digitales \cite{hadi_partovi}. Sin embargo, la enseñanza de programación a niños representa desafíos únicos, no solo debido a la complejidad del contenido, sino también por las necesidades pedagógicas específicas que conlleva. \cite{liukas_tedxcern}.

\section{Descripción del Problema}

Tal como hemos hablado en la sección anterior, enseñar programación a niños es un desafío que no podemos ignorar. Existen diversas plataformas educativas, pero muchas tienen bastantes limitaciones. Ofrecen ejercicios y evaluaciones estáticas, sin tener en cuenta una adaptación al ritmo o habilidades del estudiante. Además, la mayoría, únicamente se centran en la corrección binaria del código. Por ello, ignoran aspectos como la eficiencia del mismo y el uso de buenas prácticas de programación. Este enfoque limitado no solo afecta a los estudiantes. También deja a profesores y administradores sin las herramientas necesarias para interactuar de una forma efectiva con el sistema. Es por eso que hay una necesidad clara e importante de una solución más completa y adaptativa para la enseñanza de programación a niños.

\section{Enfoque}
Este TFG adopta un enfoque multidisciplinario, combinando elementos de pedagogía, tecnología educativa y diseño de sistemas. La naturaleza del trabajo exige la interacción de varias disciplinas con tal de abordar de manera eficaz los desafíos presentados por la enseñanza de programación a niños.

\subsection{Objetivos del Proyecto}

El propósito principal de este TFG es desarrollar un enfoque integral para la enseñanza de programación a niños, teniendo en cuenta las particularidades y desafíos que este proceso conlleva. Los objetivos específicos se detallan a continuación:

\begin{enumerate}  
    \item \textbf{Identificación y evaluación de desafíos pedagógicos:} Analizar los retos específicos asociados con la enseñanza de programación a niños, incluyendo las necesidades pedagógicas y la complejidad del contenido.
    
    \item \textbf{Diseño del Sistema de Tutorización Inteligente (ITS):} Crear una estructura para un ITS que sea adaptativo y personalizado, atendiendo a las necesidades individuales de aprendizaje de cada niño.
    
    \item \textbf{Desarrollo de una metodología adaptativa para el ITS:} Formular una metodología que permita al ITS ajustarse basándose en el progreso y las respuestas de los estudiantes, mejorando así la eficacia del proceso de aprendizaje.
    
    \item \textbf{Integración y optimización del ITS en una plataforma web:} Implementar el ITS en una plataforma web amigable y accesible, con un enfoque en la experiencia del usuario, asegurando que sea intuitiva y efectiva para niños y educadores.
    
    \item \textbf{Experimentación y análisis práctico:} Probar el sistema en contextos educativos reales, recopilando datos y feedback para evaluar su eficacia y realizar los ajustes necesarios.
    
\end{enumerate}

\subsection{Desarrollo Ágil}
Para abordar la dinámica cambiante y las necesidades de adaptabilidad del proyecto, se ha optado por usar un enfoque de desarrollo ágil para la implementación del software. Esto, favorece un proceso iterativo e incremental, permitiendo varios ajustes rápidos ante cualquier cambio. Además, de ofrecer la posibilidad de realizar entregas parciales que se pueden evaluar y ajustar en fases iniciales.

Esta estrategia de desarrollo asegura flexibilidad, fomenta la colaboración y se centra en el usuario final. Garantizando, así, que las soluciones implementadas se alinean con las necesidades del ámbito de la enseñanza y el aprendizaje.

\section{Colaboración con Codelearn S.L.}

Este proyecto no habría sido posible sin la colaboración estrecha con Codelearn S.L. \cite{codelearn}, una destacada entidad de la educación tecnológica para niños. A pesar de los avances y logros de Codelearn S.L., se identificó la necesidad de evolucionar y adaptar su sistema pedagógico para hacerlo más personalizado y efectivo. Esta colaboración ha permitido fusionar la experiencia y conocimientos de ambas partes, con el objetivo de desarrollar un ITS que permita mejorar la forma en que se enseña.

\section{Estructura de la memoria}

En esta sección se describe la estructura de esta memoria. Concretamente, se organiza en nueve capítulos principales, acompañados de anexos, que detallan el desarrollo. 

\begin{enumerate}
    \item  \textbf{Introducción} (véase capítulo \ref{chap:introduction}): Establece el contexto de la programación en la era actual, describe el problema que aborda el TFG, presenta los objetivos, el enfoque metodológico, la colaboración con Codelearn S.L., y esta misma sección que esboza la estructura del documento.
    \item  \textbf{Estado del Arte y Trabajos Relacionados} (véase capítulo \ref{chap:stateoftheart}): Profundiza en los Sistemas de Tutorización Inteligente (ITS), su evolución, aplicaciones actuales en diversos campos, limitaciones, desafíos y su papel en la educación moderna, incluyendo una comparativa de plataformas de programación.
    \item \textbf{Planificación} (véase capítulo \ref{chap:planification}): Desglosa las fases del proyecto junto con el presupuesto detallado que cubre licencias, recursos materiales, costes de personal y otros gastos.
    \item \textbf{Análisis} (véase capítulo \ref{chap:analisis}): Incluye el análisis de usuarios objetivo, requisitos educativos, modelos de casos de uso, escenarios de uso, y la seguridad y acceso al sistema.
    \item \textbf{Requisitos} (véase capítulo \ref{chap:requisitos}): Define los requisitos funcionales, no funcionales y de información del sistema, detallando aspectos generales, del sistema, de estudiantes, profesores y administradores.
    \item \textbf{Diseño e Implementación} (véase capítulo \ref{chap:analisisExperimentación}): Cubre el diseño pedagógico, la arquitectura del sistema, diseño de base de datos, interfaz de usuario y la implementación de algoritmos y módulos de aprendizaje adaptativo.
    \item  \textbf{Despliegue} (véase capítulo \ref{chap:despliegue}): Describe la creación y configuración de imágenes para el despliegue del sistema y el uso de Podman.
    \item \textbf{Pruebas y Resultados} (véase capítulo \ref{chap:resultadosExperimentales}): Presenta las pruebas realizadas y los resultados obtenidos a partir de estas.
    \item \textbf{Conclusiones y Trabajo Futuro} (véase capítulo \ref{chap:conclusiones}): Ofrece conclusiones derivadas del trabajo realizado, posibles direcciones futuras para la investigación y una reflexión personal sobre el proyecto.
\end{enumerate}    

Además, los anexos proporcionan información adicional sobre recursos informáticos (véase anexo \ref{chap:recursos}), guiones de pruebas con los niños (véase anexo \ref{chap:guion}), ficheros utilizados para el despliegue (véase anexo \ref{podmananddockerfile}), un manual de usuario (véase anexo \ref{imagenessistema}), un glosario y la bibliografía.

    % Estado del arte
    \chapter{Estado del Arte y Trabajos Relacionados} \label{chap:state_of_the_art}

En este segundo capítulo tiene como objetivo presentar todo el conocimiento
teórico adquirido en preparación para la implementación y adaptación de una
ITS, así como el estado actual de esta tecnología.

Este estudio está basado principalmente en las contribuciones presentadas en las conferencias de la \textit{International Conference on Intelligent Tutoring Systems} de los años 2016 \cite{ITS2016}, 2020 \cite{ITS2020} y 2021 \cite{ITS2021}. Estas ediciones han sido fuentes valiosas de información sobre las últimas tendencias, técnicas y avances en el campo de la tutoría inteligente. A lo largo de este capítulo, se explorará el papel de los Sistemas de Tutoría Inteligente en la educación moderna, las metodologías empleadas para adaptar el aprendizaje a las necesidades individuales de los estudiantes y las perspectivas actuales y futuras en la adaptación y personalización del aprendizaje en línea.

\section{Definición de los ITS}

\section{Evolución de los ITS}

La idea de ser capaces de utilizar la tecnología para personalizar y mejorar la experiencia de aprendizaje no es nueva. Sin embargo, con el crecimiento exponencial de las tecnologías de la información y la computación, esta idea ha ido experimentando transformaciones radicales, dando lugar a lo que se conoce hoy en día como Sistemas de Tutorización Inteligente. A continuación, se presenta una pequeña revisión de la evolución de los ITS, trazando su desarrollo desde sus inicios hasta su estado actual.

Los ITS nacieron en la década de 1970, bajo la idea de que la tecnología podía ser utilizada para simular la tutoría humana. Estos primeros sistemas, limitados por la capacidad computacional de la época, se centraban en unos conocimientos muy específicos, como podían ser la geometría o el álgebra. Utilizaban una lógica básica y simple para ofrecer retroalimentación y una guía basada en las respuestas de los estudiantes. A pesar de su simplicidad, fueron representación de un gran paso hacia la individualización del aprendizaje.

Con el progreso en la investigación de la Inteligencia Artificial (IA) en la década de 1980, los ITS comenzaron a incorporar técnicas más avanzadas. Estos sistemas no solo eran capaces de evaluar las respuestas de los estudiantes, sino también de comprender parcialmente el proceso de pensamiento detrás de esas respuestas. Así, el feedback proporcionado se volvía más personalizado y adaptativo, reflejando mejor las necesidades individuales de cada alumno.

La siguiente década trajo consigo una mayor comprensión de la importancia de la interacción humano-computadora en los procesos educativos. Donde se desarrollaron las ITS con unas interfaces más intuitivas y amigables para el usuario. Estos sistemas comenzaron a ofrecer unos mejores entornos de aprendizaje, incorporando multimedia, simulaciones y otros recursos interactivos. La idea era crear un entorno donde el estudiante no solo recibiera retroalimentación, sino que también pudiera explorar y aprender de una manera más autónoma.

En el siglo XXI, la proliferación de \textit{Big Data} y el avance en técnicas de aprendizaje profundo permitieron que los ITS alcanzaran un nuevo nivel de sofisticación. Con la capacidad de procesar grandes cantidades de datos, estos sistemas modernos pueden detectar patrones y tendencias en el comportamiento del estudiante que antes eran invisibles. Usando esta información para predecir las necesidades de aprendizaje futuras y adaptar el contenido en tiempo real, proporcionando una experiencia de aprendizaje aún más personalizada.

El potencial de los ITS sigue siendo muy amplio. Con la constante evolución de la tecnología y una comprensión cada vez más profunda de la pedagogía, es probable que veamos sistemas aún más avanzados en el futuro cercano. Incluso, se espera que la integración de realidades virtual y aumentada, junto con avances en IA, genere ITS que ofrezcan experiencias de aprendizaje inmersivas y altamente adaptativas. 

\section{Aplicaciones actuales de los ITS}

La relevancia actual de los ITS se extiende a una variedad de dominios y, con la evolución tecnológica, su precisión y adaptabilidad han alcanzado niveles sin precedentes. Aprovechando técnicas de aprendizaje automático, estos sistemas pueden reconocer patrones en el comportamiento del usuario y adaptarse en función de ello. A su vez, mediante el análisis semántico, comprenden de manera profunda las respuestas de los estudiantes, garantizando una retroalimentación más precisa y relevante. Además, algunos ITS modernos están incorporando Realidad Virtual y Aumentada, ofreciendo entornos de aprendizaje inmersivos que permiten interacciones novedosas y atractivas con el contenido. A continuación, exploraremos algunas de las aplicaciones más prominentes de estos sistemas en el contexto contemporáneo.

\subsection{Educación K-12}

Los ITS se han establecido como herramientas fundamentales para complementar el aprendizaje tradicional dentro del ámbito escolar, ya sea en la educación primaria o en la eduación secundaria. Sistemas como \textit{Dreambox} \cite{dreambox} para matemáticas o \textit{Fast Forword} \cite{fastforword} para la lectura, utilizan técnicas de tutorización inteligente para adaptar el contenido a las necesidades individuales de cada estudiante, permitiendo una educación más personalizada.

\subsection{Educación Superior}

Las universidades y colegios alrededor del mundo están empezando a integrar ITS en sus plataformas de aprendizaje en línea. Sobre todo, estas instituciones los han comenzado a usar para los cursos masivos en línea (MOOCs), donde el sistema se encarga de tutorizar a miles de estudiantes simultáneamente, sin una presencia física importante del profesor.

Por ejemplo, el sistema \textit{edX} \cite{edx}, iniciado por Harvard y el MIT, ha integrado algoritmos de tutorización inteligente para ofrecer retroalimentación personalizada a los estudiantes en sus cursos MOOCs con más inscritos. Asimismo, la Universidad de Stanford, a través de la plataforma \textit{Coursera} \cite{coursera}, también ha experimentado con herramientas de tutorización para mejorar la retención y comprensión de los estudiantes adscritos. 

\subsection{Formación y Capacitación Corporativa}

El mundo empresarial también ha reconocido el valor estos sistemas para la formación continua de sus empleados. Las grandes corporaciones han comenzado a implementarlos dentro de sus programas de formación, permitiendo la identificación y corrección de carencias de habilidades en tiempo real.

Grandes corporaciones como Google y Microsoft ya han comenzado a implementarlo en sus programas de formación interna. Por ejemplo, el \textit{Google's IT Support Certificate} \cite{googleitsupport}, un programa desarrollado para formar especialistas de soporte en TI, utiliza los ITS para guiar a los estudiantes a través de una serie de labores técnicas, identificando áreas de mejora en tiempo real.

\subsection{Salud y Medicina}

En el campo médico, también están desempeñando un papel crucial en la formación de futuros profesionales. Concretamente, \textit{Touch Surgery Simulator} \cite{touchsurgery} permite a los médicos practicar procedimientos quirúrgicos en un entorno virtual, sino que también proporciona retroalimentación detallada basada en el rendimiento y decisiones tomadas durante la simulación.

\subsection{Defensa y Aviación}

Es importante destacar que la aviación y la defensa han estado al frente en la adopción de simuladores con capacidades de tutorización. Un gran ejemplo es el \textit{F-35 Lightning II Training System} \cite{f35training}. Este ayuda a capacitar los pilotos de combate, proporcionando escenarios adaptativos basados en las habilidades y el progreso del piloto.

\subsection{Juegos Educativos y Simuladores}

Los juegos educativos, como \textit {DragonBox} \cite{dragonbox} que enseña matemáticas, o \textit {Civilization EDU} \cite{civilizationedu} que instruye sobre historia y estrategia, también integran sistemas de tutorización inteligente para adaptar el contenido y los desafíos a las habilidades del jugador, asegurando una curva de aprendizaje apropiada.


\section{Limitaciones y Desafíos}

A pesar de las prometedoras capacidades y aplicaciones de los ITS en diversos contextos educativos y de formación, estos sistemas no están exentos de limitaciones y desafíos. 

La interacción con estos sistemas no logra capturar la profundidad y empatía de un tutor humano, lo que puede generar desconexión o frustración por parte del estudiante. Estos sistemas, siendo fruto de una confluencia de disciplinas muy variadas, plantean retos en su diseño y actualización, especialmente cuando se busca mantener su relevancia en campos en constante evolución.

La adaptabilidad de estos sistemas también se encuentra en una línea delicada. Por ser lo suficientemente generales para diversos contextos. Asimismo, como específicos para brindar una retroalimentación precisa. Siendo este equilibrio esencial para garantizar su eficacia. Cabe destacar, que con la recopilación constante de datos del usuario, se debe tener especial delicadeza en cuanto a la privacidad y seguridad. Es imperativo que estos sistemas ofrezcan unas garantías robustas con tal de proteger la información del alumno.

Por otro lado, a pesar de su potencial democratizador en el acceso a la educación, los ITS corren el riesgo de convertirse en herramientas elitistas, accesibles solo para aquellos con los medios para obtenerlos. Por ello, es crucial que se trabaje en su accesibilidad y asequibilidad. Y si bien estos sistemas ofrecen métricas de rendimiento, medir su impacto real en el aprendizaje y en la transferencia práctica del conocimiento sigue siendo un gran reto.

En resumen, se podría decir que los ITS se ven como una herramienta revolucionaria, pero es fundamental enfrentar y superar los desafíos mencionados para garantizar que su implementación sea efectiva y equitativa en todos los contextos educativos.

    % Planificación del trabajo
    \input{secciones/03_ObjetivoPlanificación}

    % Desarrollo bajo sprints:
    \chapter{Análisis y decisión tecnológica} \label{chap:objetivosPlanificación}


    % Experimentación
    \chapter{Análisis y Experimentación} \label{chap:analisisExperimentación}



    % Analisis de resultados
    \chapter{Pruebas y Resultados} \label{chap:resultadosExperimentales}


     % Conclusiones
    \chapter{Conclusiones y Líneas de Trabajo Futuro} \label{chap:Conclusiones}


    \newpage
    \appendix
    \addcontentsline{toc}{chapter}{Anexos}
    \begin{appendices}

\chapter{Recursos Informáticos} \label{chap:recursos}

En este apéndice se detalla todos los recursos de hardware y software que se han usado durante la ejecución de este proyecto. Esta información es esencial para entender el entorno de desarrollo y ejecución, permitiendo así una mejor replicación o continuación del trabajo realizado.
    
\section{Hardware}

Para el desarrollo y ejecución de este proyecto, se ha utilizado el siguiente hardware:

\begin{itemize}
    \item Ordenador portátil \textit{Xiaomi TM1703}.
    \begin{itemize}
    \item Arquitectura de 64 bits.
    \item 8 GB de Memoria RAM.
    \item Procesador \textit{ Intel(R) Core(TM) 5-8250U CPU @ 1.60GHz}.
    \end{itemize}
\end{itemize}

\section{Software}

En cuanto al software, se han utilizado las siguientes herramientas y plataformas:

\begin{itemize}
    \item \textbf{Sistema Operativo}: Ubuntu 20.04.6 LTS.
    \item \textbf{Contenedores}: Docker, versión 24.0.5.
    \item \textbf{IDE}: Visual Studio Code 1.83.1
    \item \textbf{Base de datos}: PostgreSQL 16.0
    \item \textbf{Lenguaje de programación}: Python 3.8.10
    \item \textbf{Procesador de textos}: LaTex 3.14159265-2.6-1.40.20
    \item \textbf{Diseño de E-R}: draw.io
    \item \textbf{Modelado de datos}: PostgreSQL 16.0
    \item \textbf{Diagrama de flujo y arquitectura}: PlantUML v1.2023.12
    \item \textbf{Mockups}: Wireframe.cc
    \item \textbf{Bibliotecas Python}:
    \begin{itemize}
        \item Flask \cite{flask}
        \item Flask-Login \cite{flask-login}
        \item Flask-SQLAlchemy \cite{flask-sqlalchemy}
        \item Flask-CORS \cite{flask-cors}
        \item Flask-Bcrypt \cite{flask-bcrypt}
        \item python-dotenv \cite{python-dotenv}
        \item datetime \cite{datetime}
        \item werkzeug.utils \cite{werkzeug}
        \item radon.complexity \cite{radon}
    \end{itemize}
\end{itemize}
    
\chapter{Guion pruebas con niños}

Este guion tiene como objetivo guiar las pruebas de la plataforma de programación para niños. Se busca evaluar la usabilidad, el interés generado y la eficacia educativa del sistema.

\section{Preparación}
\begin{itemize}
    \item Preparar un ambiente tranquilo y sin distracciones.
    \item Asegurarse de que la plataforma funcione correctamente.
    \item Tener a mano material para tomar notas.
\end{itemize}

\section{Participantes}
\begin{itemize}
    \item Edades entre 7 y 18 años.
    \item Con y sin experiencia previa en programación.
\end{itemize}

\section{Actividades}
\subsection*{Introducción y Consentimiento}
\begin{itemize}
    \item Explicar el objetivo de la prueba.
\end{itemize}

\subsection*{Prueba de Usabilidad}
\begin{itemize}
    \item Navegar por la plataforma.
    \item Completar un ejercicio.
    \item Utilizar las herramientas de ayuda o tutoriales, si están disponibles.
\end{itemize}

\subsection*{Prueba de Interés}
\begin{itemize}
    \item Observar las reacciones de los niños mientras usan la plataforma.
    \item Preguntarles qué les gusta y qué no.
\end{itemize}

\subsection*{Prueba de Eficacia Educativa}
\begin{itemize}
    \item Evaluar el conocimiento adquirido con un pequeño cuestionario oral.
\end{itemize}

\section{Feedback y Conclusiones}
\begin{itemize}
    \item Pedir a los niños que proporcionen comentarios.
    \item Analizar los resultados y planificar mejoras para la plataforma.
\end{itemize}


\chapter{Manual de uso de la Plataforma}

\end{appendices}
    \newpage
    \section*{Glosario}

\begin{description}
    \item[ITS:] Sistema Inteligente de Tutorización (Intelligent Tutoring System). Es un sistema que utiliza tecnologías de inteligencia artificial para proporcionar enseñanza personalizada y apoyo educativo a los estudiantes, adaptándose a las necesidades individuales de aprendizaje.
    \item[Hash y Salting:] Técnicas de seguridad para el almacenamiento seguro de contraseñas. El "hashing" convierte una contraseña en una cadena de caracteres única, mientras que "salting" añade datos aleatorios al hash para aumentar la seguridad.
    \item[HTTPS:] Protocolo seguro para la transferencia de datos en la web. Utiliza el protocolo SSL/TLS para encriptar la información transmitida.
    \item[Flujo:] Secuencia de pasos o actividades en un proceso o escenario. En el contexto de programación, se refiere a la ejecución ordenada de instrucciones en un programa.
    \item[Web Responsive:] Diseño web adaptable a diferentes dispositivos. Utiliza hojas de estilo en cascada (CSS) para ajustar el diseño según el tamaño de la pantalla.
    \item[Gamificación:] Uso de elementos de juego en entornos no lúdicos. Se emplea para aumentar la motivación y la participación en actividades que normalmente se considerarían tediosas.
    \item[Aprendizaje Basado en Problemas (ABP):] Método pedagógico que utiliza problemas reales o simulados como punto de partida para el aprendizaje.
    \item[API RESTful:] Interfaz de programación de aplicaciones que sigue los principios REST para la comunicación entre sistemas.
    \item[Backend:] Parte del sistema que se encarga de la lógica de negocio, el acceso a la base de datos y la comunicación con diferentes servicios.
    \item[Clave Primaria:] Campo único en una tabla de base de datos que se utiliza para identificar registros de forma única.
    \item[Clave Foránea:] Campo en una tabla de base de datos que se utiliza para establecer una relación con otra tabla.
    \item[Normalización:] Proceso de organización de una base de datos para reducir la redundancia y mejorar la integridad de los datos.
    \item[Análisis de Árbol de Sintaxis Abstracta (AST):] Técnica de análisis de código que genera un árbol para representar su estructura sintáctica.
    \item[Complejidad Ciclomática:] Métrica que mide la complejidad de un programa basada en el número de caminos independientes a través del código.
    \item[Función Indicadora:] Función matemática que toma el valor de 1 si una condición específica se cumple y 0 en caso contrario.
    \item[Requisitos de Aprendizaje:] Objetivos educativos específicos que un estudiante debe alcanzar en un sistema de aprendizaje adaptativo.
    \item[Tasa de Fracaso:] Proporción de ejercicios o tareas que un estudiante no ha podido completar con éxito.
    \item[Normalización:] Proceso de ajustar valores medidos en diferentes escalas a una escala común.
    \item[Ponderación:] Asignación de importancia relativa a diferentes criterios o variables en un cálculo.
    \item[Análisis Estático:] Examinar el código en busca de problemas sin ejecutar el programa.
    \item[Adaptabilidad:] Capacidad de un sistema para ajustarse dinámicamente a las necesidades o comportamientos del usuario.
    \item[Fitness:] Métrica que evalúa la calidad o eficacia de una solución en comparación con una solución de referencia.
    \item[Indicadores Cuantitativos:] Métricas numéricas utilizadas para evaluar el rendimiento o eficacia de un sistema.
    \item[Usabilidad:] Medida de la eficacia, eficiencia y satisfacción con la que los usuarios pueden realizar tareas en un sistema.
    \item[Metodología de Pruebas Individuales:] Enfoque de evaluación que implica probar el sistema con cada usuario de forma separada para obtener retroalimentación más precisa.
    \item[Optimización:] Proceso de hacer cambios en un sistema para mejorarlo en términos de eficiencia, eficacia o experiencia del usuario.
\end{description}

    \addcontentsline{toc}{chapter}{Glosario}
    \printglossary
    \bibliographystyle{plain}
    \bibliography{Bibliografia/referencias}
    \addcontentsline{toc}{chapter}{Bibliografía}
    \newpage



\end{document}
